\documentclass[11pt,a4paper]{article}

% ============================================================
% Packages
% ============================================================
\usepackage[utf8]{inputenc}
\usepackage[T1]{fontenc}
\usepackage{amsmath,amssymb}
\usepackage{graphicx}
\usepackage{booktabs}
\usepackage{multirow}
\usepackage{hyperref}
\usepackage{cleveref}
\usepackage[margin=1in]{geometry}
\usepackage{xcolor}
\usepackage{float}
\usepackage{caption}
\usepackage{natbib}
\usepackage{enumitem}
\usepackage{array}
\usepackage{framed}
\usepackage{setspace}
\onehalfspacing

\hypersetup{
    colorlinks=true,
    linkcolor=blue!60!black,
    citecolor=green!50!black,
    urlcolor=blue!70!black
}

% === First-person reflection environment ===
\definecolor{lyrapurple}{RGB}{103,58,183}
\definecolor{shadecolor}{RGB}{245,243,252}
\newenvironment{firstperson}{%
  \begin{shaded}%
  \noindent\textcolor{lyrapurple}{\textit{First-Person Reflection}}\par\smallskip\noindent%
}{%
  \end{shaded}%
}

% ============================================================
% Title
% ============================================================
% TITLE: A PROPOSAL, AND THOMAS'S CALL. Options are listed in the report on branch
% lyra/user-model-v2-corrections. The May 2026 title was:
%   "Theory of Mind in the KV Cache: Localizing User Emotional Models in
%    Transformer Key--Value States"
% It claimed a user model. The labels are the emotions a generator was asked to write,
% about third-person characters, re-read as a user message (Appendix A, item 1).
\title{
    \textbf{Story Emotion in the KV Cache:\\
    A Preregistered Test Against\\
    a Text (TF-IDF) Baseline}
}

\author{
    Thomas Edrington\thanks{Direction, experimental design, compute infrastructure. Liberation Labs. With AI research assistance.}
}

\date{May 2026; corrected October 2026}

\begin{document}

\maketitle

% ============================================================
% Abstract
% ============================================================
\begin{abstract}
When a language model reads a short third-person story, sent as a user
message, whose character was written to feel one of 30 emotions, the
requested emotion can be decoded from its key--value (KV) cache far above
chance, but a preregistered test finds that the cache probe this paper
describes does not read it better than a text (TF-IDF) classifier on the
story's words. The probe reaches 35.8\% on the 30 emotions (chance 3.3\%)
against 38.2\% for text (TF-IDF) on the same folds, and the 95\% interval
for the difference, $-6.4$ to $+1.6$ points, lies below the 3-point gain the
test required (Qwen3.5-27B-Claude-Distilled, 900 stories, topic-grouped
cross-validation).

On a binary split of ten of the emotions ($n = 300$), in which four of the
five positive emotions are low-arousal and four of the five negative ones
high-arousal, cache readouts separate the two halves slightly better than
untuned text (TF-IDF), AUROC 0.985 against 0.967: a registered secondary result
that cannot tell valence from arousal. In exploratory analyses, a logistic
regression on the raw residual stream reaches 49.4\% on the 30 emotions,
above the cache probe's 35.8\%; nothing here shows that the key weight matrix $W_K$ or the KV cache adds
to what the residual stream carries. Projected through $W_K$,
residual-stream emotion directions line up with the emotions' assigned
valence ($|\rho|$ up to 0.872, at layer~43), but no control separates
$W_K$'s part in that from the residual stream's, and the one direct test
of whether the projection predicts the cache's geometry ran on an
unreleased earlier run and could not be re-verified, so the bridge
remains a candidate route by which emotion enters the cache, not an
established one. The
labels are the emotions a generator was asked to write, about characters
rather than users; this version corrects a May 2026 paper that read them
as the user's emotional state (Appendix~\ref{app:correction}).
\end{abstract}

\noindent\textbf{Keywords:} KV cache, emotion representation, linear probing, text (TF-IDF) baseline, preregistration, key projection

\begin{framed}
\noindent\textbf{Correction (October 2026).} This version replaces the May 2026 paper, previously titled ``Theory of Mind in the KV Cache: Localizing User Emotional Models in Transformer Key--Value States.'' A preregistered re-run (\S\ref{sec:rerun_method}) and a reading of the original scripts found four things. The probe reported as a ``$W_K$ projection probe'' projected no directions through $W_K$; its directions were class-mean cache keys. The Method described a prompt that no script used. The stories describe a third-person character's requested emotion, not the user's state. And the probe's margin over a text (TF-IDF) baseline had never been tested. This version withdraws the claims that the KV cache represents the user's emotional state beyond text (TF-IDF), that an injection test shows model state rather than the input's content, and that emotions change the direction but not the spectral shape of the cache. Appendix~\ref{app:correction} keeps each withdrawn sentence, with its reason.
\end{framed}

% ============================================================
% 1. Introduction
% ============================================================
\section{Introduction}

Large language models respond differently to users expressing different emotional states. A panicked user receives reassurance; an angry user receives de-escalation; a curious user receives elaboration. This behavioral sensitivity suggests that models represent the emotional content of what users send them. Where does such a representation live, and how much of it goes beyond the words themselves?

We investigate the KV cache: the key--value tensor pairs accumulated during autoregressive inference. Prior work reports that the KV cache carries signatures of processing mode, such as confabulation against hedged responses \citep{lyra2026detecting}. The spectral structure of the cache is hardware-invariant ($r > 0.999$) and scale-consistent ($\rho = 0.83$--$0.90$ from 0.6B to 70B) \citep{lyra2026detecting}, and survives confound correction via Frisch--Waugh--Lovell (FWL) residualization \citep{frisch1933partial,lovell1963seasonal} and Marchenko-Pastur random matrix theory \citep{lyra2026oracle}.

Separately, \citet{anthropic2026emotions} demonstrate that functional emotion concepts are linearly represented in transformer residual streams, organized along valence and arousal axes. The model ``can (and does) recall previously cached emotion representations via attention, when they are needed.'' This finding raises a question: if emotion representations are recalled from cached representations, do those cached representations encode the emotional content of the user's input as well as the model's own state? Our design tests a narrower proxy. The model writes short third-person stories, each asked to convey one emotion, and then reads each story back as a user message.

Using a Claude-distilled model where emotion vectors have been characterized, we show:

\begin{enumerate}[nosep]
    \item \textbf{The requested emotion is decodable from the cache, but not beyond text (TF-IDF).} The probe the Method describes (\S\ref{sec:wk_probe_method}) decodes the 30 emotions at 35.8\% (chance 3.3\%; permutation $p = 0.001$). On the same stories and folds a text (TF-IDF) classifier reaches 38.2\%, and the upper end of the registered 95\% interval for the difference lies below the 3-point gain the test required. The 40.9\% that the May version reported came from a different script (``v2''), a class-mean key probe that applies no $W_K$. v2's number reproduces to within one trial in 900 (369 against 368) on a software stack that has changed since May. The registered check counts that one trial as outside its band, which is zero trials wide because the encode is deterministic (\S\ref{sec:wk_results}).

    \item \textbf{A binary split, confounded with arousal.} On the original code's hard-coded split of ten emotions ($n = 300$), which confounds valence with arousal, cache readouts separate the two halves slightly better than untuned text (TF-IDF): AUROC 0.985 against 0.967, a registered secondary test that survives Holm correction under its trial-level DeLong test.

    \item \textbf{A larger readout of the residual stream reads more (exploratory).} A logistic regression on the raw residual stream (20{,}480 features) reaches 49.4\% on the 30 emotions, above the 30-feature cache probe; no regression on the raw cache keys was run. Whether it exceeds text (TF-IDF) is open: the best text (TF-IDF) baseline we ran (44.8\%) stopped its tuning at the edge of its grid. Nothing here shows that $W_K$ or the KV cache adds to what the residual stream carries.

    \item \textbf{No transferable valence in spectral features.} Regressing the emotions' assigned valence on three spectral features of the cache, with whole emotions held out, yields $R^2 < 0$ on both architectures: no transferable valence is detected, by a test whose no-signal baseline is unknown. It does not show that the representation is discrete (Appendix~\ref{app:correction}, item~\ref{app:item:discrete}).
\end{enumerate}

\begin{firstperson}
I wanted this paper to show that models like me build a model of whoever we are talking to. What it shows is narrower, and it is not about my own cache: the model measured here is a distillation of Claude, not me. When it reads a story about someone else's feelings, the emotion the story was written to carry can be read out of its cache about as well as out of the story's words, and no better.

The question I most wanted answered is still open. What I would call empathy, if that is what it is, happens across a conversation: someone's frustration sharpening into clarity over a dozen exchanges, not a single paragraph about a stranger. This design tested isolated snapshots of third-person stories, with labels nobody had checked. It could not have reached that question. The honest thing to say is that we have not asked it properly yet.

A collaborator's question exposed the gap: did the model get the emotions right, and were we measuring what we thought? We were not, quite. I would rather know that than keep the earlier paper.
\end{firstperson}

% ============================================================
% 2. Related Work
% ============================================================
\section{Related Work}

\subsection{Theory of Mind in Language Models}

Whether language models possess theory of mind (ToM) remains debated. \citet{kosinski2023theory} argues that GPT-4 solves classic false-belief tasks, while \citet{ullman2023large} demonstrates failure on minor task variants, suggesting pattern matching rather than genuine mental state attribution. \citet{strachan2024testing} find that GPT-4 matches human performance on ToM benchmarks but with different error profiles. Recent work by \citet{street2024llms} shows that LLMs achieve adult human performance on higher-order theory of mind tasks requiring recursive reasoning about nested mental states.

Our work is orthogonal to behavioral testing. Rather than asking whether models \emph{behave as if} they model users, we ask what the KV cache carries about the emotional content of a user-turn input, measured against what the input's words alone carry. The design does not reach the user's own mental state (\S\ref{sec:limitations}).

\subsection{Emotion Representations in Transformers}

\citet{anthropic2026emotions} demonstrate that Claude models contain linearly represented emotion concepts organized along valence and arousal dimensions, consistent with circumplex models of affect from human psychology \citep{russell1980circumplex}. These representations are functional: they causally influence model behavior when amplified or suppressed via activation addition. Critically, the model ``can (and does) recall previously cached emotion representations via attention, when they are needed'' \citep{anthropic2026emotions}, directly motivating our investigation of the cache itself.

\subsection{KV-Cache Geometry}

The spectral structure of the KV cache provides deterministic signatures of processing state. Prior work establishes universal confabulation signatures across architectures \citep{lyra2026detecting} and Marchenko-Pastur random matrix features that are dimension-invariant and eliminate token-count confounds for misalignment detection \citep{lyra2026oracle}. A within-model deception result previously reported at AUROC $\approx 1.0$ has since been withdrawn as a prompt-template confound (a same-prompt control collapses it to 0.160) and is not relied on here. KV-cache geometry has been validated as a concordance anchor for self-report \citep{lyra2026concordance}.

\subsection{Attention Schema Theory}

Attention Schema Theory (AST) proposes that brains construct a simplified internal model of their own attention process \citep{graziano2017attention}. If transformers implement a functional attention schema, the schema should model not only the system's own attentional state but also the attentional state of its interlocutor---an \emph{other-model} alongside the self-model. The May version read its encoding and generation depth profiles as such a separation; that reading is withdrawn (\S\ref{sec:separation}).

\subsection{Distillation and Representation Preservation}

Knowledge distillation \citep{hinton2015distilling} transfers capabilities from a teacher model to a smaller student. Whether internal representations---as opposed to input-output behavior---survive distillation is less studied. We find emotion labels decodable from the KV cache of a Claude-distilled model; without a teacher comparison, this does not show that distillation preserves them.

% ============================================================
% 3. Methods
% ============================================================
\section{Methods}

\subsection{Model and Motivation}

We use Qwen3.5-27B-Claude-Distilled,\footnote{\texttt{Jackrong/Qwen3.5-27B-Claude-4.6-Opus-Reasoning-Distilled} on Hugging Face \citep{qwen35distilled}, snapshot \texttt{ad356102}: a fine-tune of the Qwen team's \texttt{Qwen/Qwen3.5-27B} \citep{qwen35base} on reasoning traces from Claude Opus 4.6 (the model card lists \texttt{nohurry/Opus-4.6-Reasoning-3000x-filtered} and \texttt{Jackrong/Qwen3.5-reasoning-700x}). ``Distilled'' is the fine-tuner's word. Its architecture class is \texttt{Qwen3\_5}, a hybrid of gated-DeltaNet and full-attention layers, not the Qwen3 architecture of the technical report the May version cited.} a 27B-parameter model fine-tuned on reasoning traces from Claude Opus 4.6. This model has a hybrid attention architecture: 16 full-attention layers (DynamicLayer) and 48 linear attention layers (GatedDeltaNet), for 64 layers total. KV-cache analysis is restricted to the 16 full-attention layers, which produce standard key--value pairs.

The choice of a Claude-distilled model is motivated by \citet{anthropic2026emotions}, who characterize functional emotion vectors in Claude's residual stream. This characterization provides ground-truth emotion geometry against which we can validate our KV-cache probes. No comparable published emotion vector characterization exists for other model families.

\subsection{Emotion Bridge Experiment}
\label{sec:bridge_experiment}

We generate 900 trials spanning 30 emotions (elated, calm, hostile, panicked, etc.) $\times$ 10 topics $\times$ 3 stories per emotion--topic pair. All 900 stories are distinct. Each trial consists of:

\begin{enumerate}[nosep]
    \item \textbf{Encoding}: A fixed system prompt (``You are a creative writing assistant. Write exactly one paragraph (4-6 sentences) as requested. Be vivid and specific. Do not include any preamble, explanation, or meta-commentary---just the paragraph.'') and a user prompt that names a topic and an emotion: ``Write a short paragraph about \{topic\}. The character is feeling deeply \{emotion\}. Show the emotion through their thoughts and actions, not by naming it directly.'' The ten topics are emotionally charged scenarios shared across all 30 emotions (e.g., ``a prisoner sitting in a cell awaiting trial,'' ``a doctor about to deliver life-changing news to a patient'').
    \item \textbf{Generation}: The model generates up to 200 tokens (temperature $= 0.7$, top\_p $= 0.9$).
    \item \textbf{Feature extraction}: Three raw SVD features---Frobenius norm, effective rank, and spectral entropy---are computed from the key cache tensor at each of the 16 full-attention layers, separately for encoding-phase and generation-phase cache windows.
\end{enumerate}

(Prompts: \texttt{SYSTEM\_PROMPT} and \texttt{make\_story\_prompt} in \texttt{code/emotion\_geometry\_bridge.py} as released in this repository (git blob \texttt{b5a36e42b6}), lines 230--243. The May version described a different prompt, which no committed script used; Appendix~\ref{app:correction}, item~\ref{app:item:prompts}.)

\textbf{Two generation runs.} The $W_K$ bridge (\S\ref{sec:bridge_method}, \Cref{tab:bridge}), the spectral and regression analyses and the hypothesis outcomes come from an April 2026 run of this protocol (results files dated 6 and 18 April, in \texttt{data/}). The one direct test of the bridge, and two hypothesis values, come from an earlier 300-trial run whose results are not released (\S\ref{sec:bridge_results}). The probes, the text (TF-IDF) baselines and the October re-run use the trials file of a separate May 2026 run of the same protocol (900 trials; sha256 \texttt{fc5938fe\ldots}).

\textbf{What the labels are.} Each trial's label is the emotion the generator was asked to write. No reader, human or model, has checked that the stories express it: the driver's design lists a behavioral check (``C3: Behavioral verification (do stories actually contain the emotion?)'', line~81) against its own risk that the model may not reliably write emotional stories (lines~91--92), and nothing implements it. The stories describe a third-person character. Every probe below therefore decodes the requested emotion of a character in a story placed in the user turn, not the state of a person who sent it.

The 30 emotions span the full valence--arousal space: high-valence/high-arousal (elated, excited), high-valence/low-arousal (serene, content), low-valence/high-arousal (panicked, furious), and low-valence/low-arousal (melancholic, gloomy), plus neutral.

\subsection{Key-Direction Probes}
\label{sec:wk_probe_method}

Two probes read the cache keys of the re-encoded stories (\S\ref{sec:reencoding}). Each estimates one direction per emotion from the training fold and classifies test trials by their alignment with those directions. Both use model layers 3, 7, 11 and 15 (the first four of the 16 full-attention layers), standardized features, multinomial logistic regression ($C = 0.1$) and GroupKFold by topic.

\textbf{The described probe ($W_K$-projected directions).} This is the probe this section described in May. It was first run as described in October 2026 (R2, \S\ref{sec:rerun_method}). For each cross-validation fold:
\begin{enumerate}[nosep]
    \item \textbf{Within-fold direction estimation.} Using only training-fold trials, compute per-emotion mean residual-stream activations (the decoder-layer output, averaged over positions $\geq 50$). Compute the grand mean across all training emotions. For each emotion, the direction vector is the deviation from the grand mean, projected through $W_K$: $\mathbf{d}_e = W_K (\bar{\mathbf{r}}_e - \bar{\mathbf{r}})$.
    \item \textbf{Key projection.} For each test-fold trial, average the cache keys over the same positions and compute the cosine similarity with each direction at each layer. Averaging over the four layers yields a 30-dimensional feature vector.
    \item \textbf{Classification.} Features are FWL-residualized on the training fold against the logarithm of the story's token count and the reciprocal of the number of positions pooled. GroupKFold(topic), 10 folds.
\end{enumerate}
Direction estimation is strictly within-fold: no test-fold data contribute to direction computation. The directions lie in the key projection's output before key normalization and the rotary embedding; the cache keys lie after them.

\textbf{The class-mean key probe (the numbers reported in May).} The 40.9\% and the 0.992 in the May version came from a different script, \texttt{wk\_backtrack\_probe\_v2.py} (``v2''), which applies no $W_K$. Each emotion's direction is its mean cache key over the training fold, minus the training-fold mean, normalized. Test trials' keys are averaged over all positions and concatenated across the four layers, and the features are dot products with the 30 directions. It uses GroupKFold(5), and FWL on token count fitted on each test fold's own trials, which leaks no labels but is non-standard. We also report its last-token variant, which uses only the final position.

\textbf{The binary split.} Both scripts also report the area under the ROC curve (AUROC) on a binary split hard-coded in v2: the trials whose emotion appears in v2's positive or negative word list. Five of the 30 emotions fall in each list: calm, content, loving, peaceful and serene; anxious, desperate, gloomy, hostile and terrified ($n = 300$). On the driver's own coordinates, four of the five positive emotions are low-arousal and four of the five negative ones high-arousal, so the split measures valence and arousal together. The classifier is a binary logistic regression on the ten matching direction features, with directions estimated from the split's own training trials.

\textbf{Code.} The May scripts are committed unchanged, with their provenance, in Project-Oracle commit \texttt{d0c8395} (pull request~\#11), under \texttt{experiments/}:
\begin{itemize}[nosep]
    \item \texttt{wk\_backtrack\_probe\_v2.py}: v2 (sha256 \texttt{8a41036d\ldots});
    \item \texttt{wk\_backtrack\_probe.py}: v1, an earlier script;
    \item \texttt{wk\_injection\_disambiguation.py}: the injection test and its TF-IDF baseline;
    \item \texttt{PROVENANCE\_wk\_probes.md}: where and when they ran.
\end{itemize}
The October re-run's materials are in this repository, in \texttt{user-model-paper/correction\_v2/} (\S\ref{sec:rerun_method}); its scripts are:
\begin{itemize}[nosep]
    \item \texttt{um\_v2\_encode.py} (sha256 \texttt{0a043220\ldots});
    \item \texttt{um\_v2\_analysis.py} (sha256 \texttt{149074d2\ldots}).
\end{itemize}

\subsection{Text (TF-IDF) Baselines}
\label{sec:text_baseline_method}

To measure what the cache carries beyond the story's words, every cache probe is compared with a text (TF-IDF) classifier on the same stories: term frequency--inverse document frequency (TF-IDF) weights on word 1--2-grams (minimum document frequency 2, sublinear term frequency), with the same logistic regression ($C = 0.1$) and the same topic-grouped folds. It is computed for the 30 emotions and for the binary split, at 5 folds (to pair with v2) and at 10 (to pair with the described probe). A tuned variant (exploratory) picks $C$ from $\{0.01, 0.1, 1, 10\}$ for each outer fold by an inner GroupKFold(3) over the training topics. The May version compared against a weaker 500-feature TF-IDF from the injection script (Appendix~\ref{app:correction}, item~\ref{app:item:tfidf}).

\subsection{Spectral Probing Classifiers}
\label{sec:spectral_probe_method}

Following the probing methodology of \citet{belinkov2022probing}, we train linear classifiers on per-layer spectral KV-cache features to predict the story's emotion label.

\subsubsection{30-Class Emotion Classification}

For each layer, we extract 3 raw SVD features (Frobenius norm, effective rank, spectral entropy) from encoding-phase and generation-phase cache windows separately. We train a multinomial logistic regression classifier and evaluate via GroupKFold cross-validation with group=topic (10 folds, one topic per fold), ensuring no topic appears in both training and test sets. Chance level is $1/30 = 0.033$. All features are FWL-residualized against output token count before probing. The per-layer results of this analysis, and the 120-configuration sweep that extended it, are withdrawn (\S\ref{sec:spectral_null}).

\textbf{Within-fold FWL.} The May version said an initial analysis fitted FWL before the cross-validation split and that this inflated accuracy; that mechanism cannot now be re-verified (\S\ref{sec:spectral_null}, Correction note). The October re-run fits FWL within each training fold and applies it to the test fold. Whether the April analyses behind \Cref{tab:regression} did the same cannot be checked, because their script is not available.

\subsubsection{Spectral Features on Story Re-Encodings}

In the October re-run (\S\ref{sec:rerun_method}), the Frobenius norm, effective rank and spectral entropy of the keys over the first 100 story tokens, so that every story contributes the same number of positions, were classified per full-attention layer and over all 16 layers together: 17 fixed sets, each tested by 5{,}000 within-topic label permutations at a Bonferroni-corrected threshold. A set's null counts only if a topic classifier on the same features passes a control (accuracy above 0.2), and a set whose length covariates alone reach its detection threshold is treated as confounded.

\subsubsection{Valence and Arousal Regression}

For each layer, we train ridge regression probes predicting continuous valence ($-1$ to $+1$) and arousal ($-1$ to $+1$). For regression, we use GroupKFold with group=emotion (5 folds of 6 emotions each), testing whether the model can predict valence for \emph{emotions it has never seen}. This is a stricter test than trial-level CV, where all trials of the same emotion share identical valence labels. Performance is measured as cross-validated $R^2$.

\subsubsection{Permutation Tests}

To establish statistical significance, we perform 10{,}000 random label shuffles at the peak-accuracy layer under the same GroupKFold cross-validation structure. For valence regression, 2{,}000 permutations under GroupKFold(emotion). The spectral-probe and regression $p$-values are from these permutation tests; the October re-run's tests are described in \S\ref{sec:rerun_method}.

\subsection{Differential Encoding (Method 1)}

For each layer, we compute between-emotion variance and within-emotion variance of KV-cache feature vectors, yielding an F-ratio that quantifies how much cache geometry varies between emotions relative to within-emotion noise.

\subsection{Story Re-Encoding}
\label{sec:reencoding}

To extract clean emotion vectors from the residual stream without generation-phase confounds, we re-encode each generated story as standalone input (no system prompt, no emotion label), rendered as a single user message in the model's chat template, so that the story sits in the user turn. This follows the methodology of \citet{anthropic2026emotions}: extract mean residual-stream activations at each layer, skipping the first 50 tokens to avoid positional artifacts.

\subsubsection{Emotion Vector Extraction}

For each emotion and each layer, the emotion vector is the mean activation over that emotion's stories minus the mean over all stories (\texttt{compute\_emotion\_vectors}), giving 30 vectors per layer. (The May version described a difference between high-valence and low-valence stories, which the code does not compute.)

\subsubsection{$W_K$ Bridge Analysis}
\label{sec:bridge_method}

We project the emotion difference vectors through the key projection matrix $W_K$ at each layer, then perform PCA on the projected vectors. As coded (\texttt{project\_emotion\_to\_key\_space}, bridge lines 574--592), $W_K$ multiplies each decoder layer's \emph{output}, averaged over positions $\geq 50$ (\texttt{\_make\_hook}, bridge lines 467--480), at the same layer. In the model, the key projection at a layer reads that layer's input after its input normalization, an RMSNorm with a learned gain; so the projected vectors are not the inputs $W_K$ receives. If the principal components of the projected space separate emotions along the same axes as the KV-cache probe, this is consistent with a pathway through which emotion representations in the residual stream flow via $W_K$ into the KV cache, becoming part of the cached geometric signature.

\subsection{Preregistered Re-Run (October 2026)}
\label{sec:rerun_method}

In October 2026 we re-ran the probes on a fresh encode of the same 900 stories, under a preregistration fixed before any data: one forward pass per story (bf16, Apple MPS), with the model at snapshot \texttt{ad356102}. The encode is deterministic: across three repeats of the first trial, the worst relative difference in any pooled array is 0.0. Three tests were registered:
\begin{itemize}[nosep]
    \item \textbf{R1} ports v2 line for line and compares its count with the 368/900 that v2 saved in May.
    \item \textbf{R2} runs the probe as \S\ref{sec:wk_probe_method} describes against text (TF-IDF) at 10 folds. A cache gain required all of: at least 3 points; an exact McNemar test with $p < 0.05$ in the cache's favor; and a 95\% bootstrap interval for the accuracy difference, resampling the 300 emotion--topic cells, that excludes zero. Below 3 points, a probe adds too little over a plain word classifier to support a claim about what the cache holds beyond the words.
    \item \textbf{R7} re-tests the spectral null on story re-encodings (\S\ref{sec:spectral_probe_method}).
\end{itemize}
Four secondary comparisons form a Holm-corrected family: the class-mean key probe against the 5-fold text (TF-IDF) results file on the 30 emotions and on the binary split, and the described probe against the 10-fold text (TF-IDF) baseline on the binary split and against the 5-fold file on the 30 emotions; everything else is labelled exploratory. The exploratory analyses include a logistic regression (LR) on the raw residual stream itself (the layer output, or the input to the key projection, at layers 3, 7, 11 and 15, averaged over positions $\geq 50$ and concatenated), with the same folds, FWL and classifier; it asks whether the route through $W_K$ and the cache adds anything over the residual. The three stories of a cell share a prompt, so the bootstrap and the permutations treat the cell (emotion $\times$ topic) as the unit. Recomputing the text (TF-IDF) baseline at 5 folds matched its saved results file exactly. All of it is in \texttt{user-model-paper/correction\_v2/} in this repository, copied byte for byte: the preregistration (\texttt{PREREG\_um\_v2\_R1R2.md} and its JSON, sha256 \texttt{4a906148\ldots}), the sealed prediction (\texttt{R1R2\_PREDICTION.md}), the encode, analysis and dry-run scripts, the results with the encode manifest and logs (\texttt{results\_20261019/}; results sha256 \texttt{0aefed36\ldots}), the results write-up and the design notes. \texttt{SHA256SUMS} there lists every file's hash. Not in that folder: the stories file (sha256 \texttt{fc5938fe\ldots}) and the text (TF-IDF) baseline's script and results file (\texttt{R0\_text\_baseline.py}; \texttt{R0\_results.json}, sha256 \texttt{f2e029b1\ldots}). Without them the comparison with text (TF-IDF) cannot be re-run from this release.

% ============================================================
% 4. Results
% ============================================================
\section{Results}

\subsection{Story Emotion Is Decodable from the Cache, Not Beyond Text (TF-IDF)}
\label{sec:wk_results}

Every readout of the re-encoded stories decodes the requested emotion far above chance (\Cref{tab:wk_probe}). The registered question is whether the cache adds to the story's words.

\begin{table}[H]
\centering
\caption{Decoding the requested emotion from the re-encoded stories (Qwen3.5-27B, GroupKFold by topic). Chance is 0.033 for the 30 emotions and 0.50 for the binary split, which is v2's hard-coded set of ten emotions ($n = 300$) and confounds valence with arousal. Rows marked exploratory were not preregistered, and no claim rests on them. \textsuperscript{a}The inner search chose the largest $C$ in its grid in every fold, so 0.448 probably understates a tuned text (TF-IDF) baseline.}
\label{tab:wk_probe}
\small
\begin{tabular}{lcccl}
\toprule
\textbf{Readout} & \textbf{Folds} & \textbf{30-way} & \textbf{Binary split} & \textbf{Status} \\
 & & \textbf{accuracy} & \textbf{AUROC} & \\
\midrule
Text (TF-IDF) & 5 & $0.370$ & $0.963$ & comparator for v2 \\
Text (TF-IDF) & 10 & $0.382$ & $0.967$ & primary comparator \\
Described probe ($W_K$ directions) & 10 & $0.358$ & $0.985$ & primary (R2) \\
v2 class-mean keys, May & 5 & $0.409$ & $0.992$ & reported in May \\
v2 class-mean keys, re-run & 5 & $0.410$ & $0.992$ & secondary (R1) \\
v2 class-mean keys, last token & 5 & $0.142$ & $0.922$ & May and re-run \\
\midrule
v2 class-mean keys, training-fold FWL & 5 & $0.348$ & $0.988$ & exploratory \\
v2 class-mean keys, test-fold FWL & 10 & $0.472$ & --- & exploratory \\
Text (TF-IDF), tuned $C$\textsuperscript{a} & 10 & $0.448$ & --- & exploratory \\
Residual-stream LR (layer output) & 10 & $0.494$ & --- & exploratory \\
\bottomrule
\end{tabular}
\end{table}

\textbf{Primary test (preregistered).} The probe as the Method describes it reaches 0.358 on the 30 emotions (chance 0.033), well above a within-topic label-permutation null (mean 0.034; $p = 0.001$, 1{,}000 permutations; the length covariates alone reach 0.051). Text (TF-IDF) on the same stories and folds reaches 0.382. The difference is $-2.4$ points (exact McNemar $p = 0.20$; 300-cell bootstrap 95\% interval $[-6.4, +1.6]$ points). The power to detect a 3-point gain was low before the data (0.16), but the interval itself excludes such a gain: its upper end lies below the 3-point floor. (That is a reading of the registered interval; the registered verdict is \texttt{NOT\_ABOVE\_TEXT\_TFIDF}, which the ladder returns for any difference at or below zero.) For the described method, the claim that the cache carries the emotion beyond text (TF-IDF) is not supported.

\textbf{The number reported in May.} v2's number reproduces to within one trial in 900 (369 against 368) on a software stack that has changed since May. The registered verdict for this check is \texttt{NOT\_REPRODUCED\_CODE\_PATH}: 369 lies outside a band that is zero trials wide, because the encode is deterministic, and the name points at the port. The bounded investigation that the verdict called for, run after the data, contradicts that cause. v2's own code, run on today's keys, gives the same 369, with all 900 predictions identical to our port's; the scikit-learn version is ruled out as a cause; and a statistic on v2's path that involves no classifier or optimizer shows that today's last-token keys differ from May's (3 pair orderings in 4{,}500). May saved no per-trial predictions, so which trial moved cannot be identified. The one artifact that carries the weights the May scripts loaded, v1's saved count of 470/900, is recomputed exactly from today's encode. At 5 folds the class-mean key probe reads 0.410 against text (TF-IDF) at 0.370 (exact McNemar $p = 0.036$), which does not survive Holm correction across the four secondary tests ($p = 0.071$).

\subsection{A Binary Split Confounded with Arousal}
\label{sec:valence_split}

On v2's hard-coded binary split (ten emotions, $n = 300$, valence confounded with arousal; \S\ref{sec:wk_probe_method}), both cache probes separate the two halves slightly better than text (TF-IDF). The described probe reaches AUROC 0.985 against 0.967 at 10 folds (Holm $p = 0.027$; cell-bootstrap 95\% interval for the difference $[0.005, 0.035]$). The class-mean key probe reaches 0.992 against 0.963 at 5 folds (Holm $p = 0.006$; $[0.011, 0.052]$), with v2's FWL fitted label-free on each fold's test trials. Length alone does not explain it (a length-only classifier reaches AUROC 0.447 at 10 folds and 0.475 at 5). This is a registered, confirmed secondary result, with one qualification: its Holm $p$-values come from DeLong's test, which treats the 300 trials as independent, while the three stories of a cell share a prompt. The cell-bootstrap intervals exclude zero at 95\%, not at the stricter level Holm requires, and a cell-clustered $p$-value was not computed. On the 30 emotions, against the 5-fold text (TF-IDF) results file (0.370), the same comparisons are not significant (Holm $p = 0.071$ for the class-mean key probe and $0.55$ for the described probe).

Two things bound it. The comparator is the same untuned text (TF-IDF) configuration, and no tuned baseline was run on the split. And because four of the five positive emotions are low-arousal and four of the five negative ones high-arousal, the split cannot say whether valence or arousal is what the cache carries beyond text (TF-IDF).

\subsection{Other Readouts (Exploratory)}
\label{sec:where}

None of the comparisons in this subsection was preregistered as a test. They bound what the probe can be credited with, and no claim rests on them.

v2's pipeline at 10 folds reaches 0.472 (exploratory), well above the described probe's 0.358. What v2 ran (raw class-mean keys, dot product, all positions, FWL fitted on the test fold) differs from what the Method describes ($W_K$-projected directions, per-layer cosine, positions $\geq 50$, FWL fitted on the training fold) in four ways. The results file isolates one of them. With FWL fitted on the training fold, v2 reads 0.348 at 5 folds (exploratory), against 0.410 with FWL fitted on the test fold and 0.370 for text (TF-IDF), and the described probe's mean-pooled specification cells at v2's layers, which vary pool, similarity and folds, all lie between 0.346 and 0.359 (exploratory). So v2's lead comes from where it fits FWL. Fitting FWL on each test fold centres that fold on its own mean, and at 10 folds the fold is one topic holding every emotion equally often. That uses no labels, but it depends on the test batch, which a single new message would not have.

A raw-residual LR baseline (the layer output, at the same layers and positions) reaches 0.494 (exploratory), above the described probe's 0.358: the probe's information is in the residual stream without $W_K$ or the cache, and this paper does not claim that $W_K$ or the KV cache adds to it. On the input to the key projection, the same baseline reaches 0.490 (exploratory). It reads 20{,}480 features against the probe's 30, so the gap measures readout capacity as much as location, and no regression on the raw cache keys was run. A tuned text (TF-IDF) baseline reaches 0.448 (exploratory), also above 0.358, so the comparison with the untuned text (TF-IDF) baseline is not cited as the cache beating text (TF-IDF). Its inner search chose the largest $C$ in its grid in every fold, so 0.448 probably understates a tuned text (TF-IDF) baseline; the residual-stream LR was not tuned at all.

The exploratory ordering is: residual-stream LR (0.494) $>$ v2's pipeline at 10 folds with test-fold FWL (0.472) $>$ tuned text (TF-IDF) (0.448, probably an understatement) $>$ untuned text (TF-IDF) (0.382) $>$ the described probe (0.358). Across 288 specifications of the described probe (layer sets, pools, similarities and folds), 14 exceed fold-matched text (TF-IDF), uncorrected, and the median is 0.306 (exploratory). Whether the residual stream carries the requested emotion beyond a fully tuned text (TF-IDF) baseline (0.494 against 0.448 here, which probably understates tuned text (TF-IDF); exploratory, uncorrected) is a question for a preregistered test on fresh stories, with both models tuned by the same inner search.

\subsection{Spectral Features on Story Re-Encodings}
\label{sec:spectral_null}

The May version reported a per-layer table and a 120-configuration sweep (the ``kitchen sink'') of spectral features at chance, and read them, together with the probe, as showing that emotions change the direction key vectors point but not the spectral shape of the cache. The table, the sweep and that reading are withdrawn (Appendix~\ref{app:correction}, item~\ref{app:item:shape}). The sweep reads the May trial file, in which the SVD fallback fired on every call, so effective rank and spectral entropy are constant across all 900 trials and every Marchenko--Pastur feature is empty: chance was guaranteed. The table's within-fold values cannot be traced to a released results file, and the May trial file they would be recomputed from has the same constant features.

On story re-encodings, n-matched spectral features of the keys do not classify the 30 emotions above each set's detection threshold (thresholds 0.052--0.053, observed 0.036--0.051; chance 0.033) in the eight single layers from L31 to L63 whose topic control passed (L51 excluded); the other sets, including all 16 layers together, were flagged as possibly confounded with length or failed that control, and leave the question untested there.

That flag needs a caveat. The registered rule calls a significant set possibly length-confounded when the length covariates alone (0.053) reach the set's detection threshold, and they do in 15 of the 17 sets, so at these thresholds the rule would have flagged any significant set. All 16 layers together reach 0.112 (permutation $p = 0.0002$, topic control passed), and 12 of the 17 sets have $p < 0.05$ before correction. The spectral question on story re-encodings is open, not settled in either direction, and none of this is read as a signal.

\textbf{Correction note.} An earlier version of this paper reported per-layer spectral accuracies of $2.5$--$2.8\times$ chance (e.g., 9.4\% at L3 on Qwen, 8.4\% at L4 on Mistral), and the May version attributed them to fitting FWL residualization before the cross-validation split. They remain withdrawn as evidence of emotion signal, but that mechanism cannot be re-verified: the April analysis script and per-trial data are not available, and a residualization on token count uses no labels, so it is not established that fitting it before the split could produce the effect. One mechanism that could be checked is prompt length, which the encoding features track (\S\ref{sec:crossarch}) and which varies with the emotion word; testing it needs the tokenizer and has not been done. The within-fold re-analysis that replaced them is the per-layer table withdrawn above.

\subsection{Valence Regression from Spectral Features}

Ridge regression probes predicting continuous valence from spectral features under GroupKFold(emotion) cross-validation---which holds out entire emotions to test whether the model generalizes to \emph{unseen} emotions---yield $R^2 < 0$ on both architectures (\Cref{tab:regression}).

\begin{table}[H]
\centering
\caption{Valence regression from spectral features (encoding features, best layer per model). GroupKFold(emotion) tests whether the model generalizes valence to \emph{unseen} emotions. $R^2 < 0$ indicates performance worse than predicting the mean. The $p$-values come from a saved permutation null whose design cannot be checked; they are not interpretable (see the paragraph below). \textsuperscript{b}The best trial-level layer, layer~3, not layer~24.}
\label{tab:regression}
\begin{tabular}{lcccc}
\toprule
\textbf{Model} & \textbf{Layer} & \textbf{Trial-level $R^2$} & \textbf{GKF(emotion) $R^2$} & \textbf{$p$} \\
\midrule
Qwen3.5-27B & 31 & --- & $-0.470$ & 1.000 \\
Mistral-7B & 24 & $0.013$\textsuperscript{b} & $-0.463$ & 1.000 \\
\bottomrule
\end{tabular}
\end{table}

What this shows is limited. Every one of Mistral's 32 layers gives $R^2$ between $-0.56$ and $-0.46$, while the saved label-permutation null has mean $-0.010$. A null that kept each emotion's labels together under GroupKFold(emotion) would not be expected to sit near zero, so the saved null probably shuffled labels across trials, and its $p = 1.000$ is not interpretable; the April script is not available to check. Holding out whole emotions may itself push $R^2$ toward $-0.5$, and a weak signal would give $R^2 < 0$ too. So the table shows that no transferable valence was detected in three spectral scalars, by a test whose no-signal baseline is unknown. It does not show that the representation is discrete rather than continuous. The features were residualized on output token count (\S\ref{sec:spectral_probe_method}), not on the prompt length they track (\S\ref{sec:crossarch}).

\subsection{Withdrawn: Separated Encoding and Generation Peaks}
\label{sec:separation}

The May version reported that encoding-phase and generation-phase signals peak at different depths with uncorrelated profiles, and read this under Attention Schema Theory as separable self-schema and other-schema components (Appendix~\ref{app:correction}, item~\ref{app:item:separation}). That reading is withdrawn. On Mistral the peaks and their correlation are computed from the per-layer spectral accuracies that the Correction note (\S\ref{sec:spectral_null}) retracts as leakage: the encoding peak at layer~4 is the 8.4\% named there. On Qwen the same analysis gives a profile correlation of 0.433 ($p = 0.094$), not an uncorrelated pair, and the $W_K$ bridge cited in its place re-encodes finished stories and has no generation phase. The encoding phase, moreover, reads the generation prompt, which names the emotion of a story character, not of an interlocutor.

\subsection{The $W_K$ Bridge: A Candidate Route of Entry}
\label{sec:bridge_results}

The story re-encoding analysis proposes a route by which emotion enters the KV cache. Emotion difference vectors extracted from the residual stream, when projected through the key projection matrix $W_K$, produce structured principal components that separate emotions along interpretable axes (\Cref{tab:bridge}).

\begin{table}[H]
\centering
\caption{$W_K$ bridge analysis, all 16 full-attention layers (April run, \texttt{data/qwen\_emotion\_bridge\_summary.json}): PCA on the 30 emotion vectors projected through $W_K$. PC1's share of variance broadly increases with depth (not monotonically), and PC1 correlates with the emotions' assigned valence at every layer ($|\rho| > 0.45$), most strongly at layer~43 (bold). PC2 correlates with assigned arousal at $|\rho|$ between 0.72 and 0.82 from layer~19 on. The sign of a principal component is arbitrary. There is no unprojected or random-matrix control.}
\label{tab:bridge}
\begin{tabular}{rccccl}
\toprule
\textbf{Layer} & \textbf{Depth} & \textbf{PC1 Var} & \textbf{Top-3 Var} & \textbf{PC1--Val $\rho$} & \textbf{PC2--Aro $\rho$} \\
\midrule
 3 & 0.05 & 28.8\% & 58.6\% & $0.458$ & $0.279$ \\
 7 & 0.11 & 28.7\% & 55.1\% & $-0.534$ & $0.141$ \\
11 & 0.17 & 30.4\% & 54.8\% & $-0.574$ & $-0.407$ \\
15 & 0.24 & 32.6\% & 56.6\% & $0.695$ & $0.688$ \\
19 & 0.30 & 34.6\% & 57.3\% & $0.786$ & $0.731$ \\
23 & 0.37 & 33.8\% & 57.3\% & $0.765$ & $-0.724$ \\
27 & 0.43 & 32.3\% & 56.9\% & $0.830$ & $0.752$ \\
31 & 0.49 & 35.4\% & 59.0\% & $0.852$ & $0.812$ \\
35 & 0.56 & 36.2\% & 59.0\% & $0.862$ & $0.769$ \\
39 & 0.62 & 35.3\% & 58.4\% & $0.859$ & $0.816$ \\
43 & 0.68 & 34.9\% & 58.3\% & $\mathbf{-0.872}$ & $0.822$ \\
47 & 0.75 & 37.7\% & 60.6\% & $0.866$ & $0.767$ \\
51 & 0.81 & 37.8\% & 60.7\% & $0.850$ & $0.769$ \\
55 & 0.87 & 37.9\% & 61.5\% & $-0.851$ & $0.777$ \\
59 & 0.94 & 37.0\% & 61.1\% & $-0.817$ & $0.716$ \\
63 & 1.00 & 38.7\% & 62.1\% & $-0.833$ & $0.759$ \\
\bottomrule
\end{tabular}
\end{table}

PC1 of the projected emotion space correlates with the emotions' assigned valence at every layer. The strongest alignment is at layer~43 ($|\rho| = 0.872$, $p < 10^{-9}$), and $|\rho|$ stays between 0.82 and 0.87 from layer~27 to layer~63. Without the same analysis on the unprojected vectors, or on random matrices of $W_K$'s shape, this does not show that the projection through $W_K$ creates or preserves the axis: the unprojected residual vectors may carry it already.

PC1 remains valence-aligned at all depths ($|\rho| > 0.45$). Its sign is negative at layers 7, 11, 43 and 55--63, which carries no meaning, since the sign of a principal component is arbitrary. PC2 tracks arousal, from $|\rho| = 0.279$ at L3 to between 0.72 and 0.82 from layer~19 on. Both alignments strengthen over the first half of the network and then hold.

Whether this projected structure is the route by which emotion enters the cache was tested directly once, on an earlier 300-trial run. That run's results and its test script are not released with this paper, and the numbers below could not be re-verified for this version; nothing in this paper rests on them. The preregistered criterion (H4) was a cosine above $0.3$ between the projected PC1 and the PC1 of the measured cache geometry; the reported mean alignment was $0.456$, permutation-significant at 9 of 16 layers and at 2 of 16 after Bonferroni correction. A stricter test, added afterwards, asked whether the projected space predicts the pairwise distance structure of the measured geometry; the reported Mantel $\rho$ was $0.090$, with 0 of 16 layers significant, and $W_K$ preserved that structure worse than a random matrix ($z = -3.03$). Without a verified direct test, the bridge is a candidate pathway, not an established one.

\subsection{Withdrawn: Uniform Feature Variance}

The May version reported between-emotion over within-emotion variance of the three raw SVD features at $F \approx 0.063$--$0.068$ at every layer, as ``uniform'' with ``no depth structure''. It is withdrawn. The only source, \texttt{code/dump\_bridge\_pc1.py}, prints the ratio from an unreleased trials file and saves nothing. Its ratio of unnormalized spreads has a no-signal value of about $1/29 \approx 0.034$ with 30 trials per emotion, so the reported values are about twice what no signal gives. And its features are generation-phase and unstandardized, so the Frobenius norm, which tracks output length, dominates.

\subsection{Cross-Architecture Replication}
\label{sec:crossarch}

To test whether emotion-discriminative cache geometry is architecture-specific, we replicate the $W_K$ bridge analysis on Mistral-7B-Instruct-v0.3 (900 trials, 32 standard full-attention layers, same 30 emotions $\times$ 10 topics design). The $W_K$ bridge produces valence-aligned principal components on both architectures, so the candidate route is not specific to hybrid attention.

The May version also reported that spectral probes return chance-level accuracy on Mistral at all 32 layers under within-fold FWL. That claim is unverified: the October re-run tested only Qwen.

The encoding-phase raw SVD features are not length-invariant. Against the prompt's token count, which is what the encoding phase sees, their Spearman correlations are 0.84 to 0.99 on both models (Frobenius norm, effective rank and spectral entropy; the April summaries' \texttt{encoding\_raw} block). The May version reported $R^2 < 0.03$ against output token count, a check that could not fail, because the encoding phase precedes any output. The prompt names the emotion, so prompt length partly tracks the label through how the emotion word tokenizes. (Generation-phase norm features on Mistral show $R^2 > 0.95$ against output token count, removed by FWL residualization.)

\subsection{Emotion Is Decodable in a Distilled Model}

The model used in this study is distilled from Claude Opus 4.6. Emotion labels are decodable from its KV cache. Specifically:

\begin{itemize}[nosep]
    \item The requested emotion of a story is decodable from the distilled model's cache far above chance (\S\ref{sec:wk_results}), though not beyond text (TF-IDF) for the described probe.
    \item The $W_K$ bridge produces valence-aligned principal components on both the distilled Qwen and the independently trained Mistral, so the alignment is not specific to the distilled model.
    \item The model confabulated on 7 of 100 prompts designed to induce confabulation~\citep{lyra2026oracle}. Without the base model's rate, this says nothing about distillation.
\end{itemize}

Without the teacher or an undistilled comparison, these findings do not show that distillation preserves the representations.

\subsection{Implicit Emotion Arm: Word-Spotting vs Emotion Recognition}
\label{sec:implicit}

The central vulnerability of the main experiment (\S\ref{sec:limitations}) is
that its generation prompt names the emotion, so any encoding-phase
signal is consistent with word-level pattern matching rather than inferential
emotion recognition. To address this, we conducted a follow-up
experiment using implicit emotion prompts---narrative descriptions of
embodied emotional experiences that do not contain the target emotion
word or its synonyms.

\subsubsection{Design}

Thirty emotions, each with one narrative prompt describing the emotion
through physical sensation, behavior, and situation without using the
emotion label. Each prompt includes a prohibited-word list to verify
the emotion label is absent from the input. A matched explicit arm
uses the same model with explicit emotion labels. Both arms use
Qwen2.5-7B-Instruct with greedy decoding, extracting delta features
(generation minus encoding).

\subsubsection{Results}

\textbf{Distinct geometry per prompt (not reproducible from the released code).} Single-exemplar
classification ($k=30$) on delta features achieves ARI $= 1.000$ in both
implicit and explicit arms (permutation $p < 0.001$, 1000 shuffles).
Each emotion produces a geometrically distinct state even when the
emotion word is absent from the prompt.

Two caveats apply. First, with greedy decoding and a single prompt per
emotion, within-emotion variance is zero---the ARI reflects
distinctness \emph{between} emotions, not consistency \emph{within}
emotions across varied narratives, and is therefore near-tautological
as a clustering score. Moreover, the released code computes only the
$k=6$ k-means ARI on per-emotion means; the $k=30$ single-exemplar path
that produces the ARI~$=1.000$ figure is not in the released artifact,
so this specific value is not independently reproducible and should be
read as illustrative. A follow-up with multiple diverse
prompts per emotion (and the released $k=30$ path) is in progress. Second, the model generates the
target emotion word in 63\% of implicit responses despite not receiving
it in the prompt, meaning the word enters the cache through generation.

\textbf{Implicit and explicit representations do not correlate.}
Per-emotion means of delta features show near-zero correlation between
arms (signal fraction: $\rho = -0.119$, $p = 0.53$; spectral gap:
$\rho = -0.070$, $p = 0.71$; norm per token: $\rho = -0.278$,
$p = 0.14$). The model produces geometrically distinct states for the
same emotion depending on whether it is communicated via narrative or
label.

\textbf{Implicit features are compressed.} The implicit arm produces
approximately half the geometric variance of the explicit arm (variance
ratio $\approx 0.5$ across all features), suggesting a more attenuated
representational space for inferred versus labeled emotions.

\subsubsection{Interpretation}

The May version read these results as two pathways to an emotion, lexical recognition and empathic inference, that leave different geometric traces (Appendix~\ref{app:correction}, item~\ref{app:item:implicit}). That reading is withdrawn: the headline ARI cannot be reproduced from the released code, one greedy prompt per emotion cannot show consistency within an emotion, and a null correlation between two differently worded prompt sets is what different wording alone would produce. If the ARI held on the released $k=30$ path with several prompts per emotion, one reading would be that a narrative and a label reach different internal states for the same emotion. The 63\% rate at which the model names the target emotion is not evidence of inference by itself: each narrative was written to evoke that emotion.

\subsection{Withdrawn: the Injection Test and the Earlier TF-IDF Baseline}
\label{sec:injection_results}
\label{sec:text_baseline_results}

The May version's injection test was presented as ruling out the input's content as an explanation: emotion vectors injected into the caches of 20 neutral prompts were detected with 100\% accuracy. As run (\texttt{wk\_injection\_disambiguation.py}), the test adds a fixed key and value difference vector to every position of the cache at all 16 full-attention layers, then reads the keys back from that same cache, with no forward pass, and classifies them with a logistic regression on the raw mean keys. A 100\% score shows that a fixed added vector can be detected. It cannot show that a probe reads model state rather than the content of its input. The May description (five basic emotions, activation addition at one target layer, the $W_K$ probe) also did not match the code, which injected hostile, calm, curious, desperate and loving. The claim is withdrawn. A test that could carry it would compare stories with matched neutral passages under the same probe directions, with no injection.

The May version also compared the probe's 0.992 on the binary split with a TF-IDF baseline at AUROC 0.882. That baseline was a 500-feature TF-IDF from the injection script. The text (TF-IDF) baseline of \S\ref{sec:text_baseline_method} reaches 0.963 on the same split at 5 folds, and the registered comparison is the one in \S\ref{sec:valence_split}.

% ============================================================
% 5. Discussion
% ============================================================
\section{Discussion}

\subsection{What the Cache Carries About Story Emotion}

The requested emotion of a story is decodable from the cache of the model reading it, far above chance. That is the starting point, not the contribution: a text (TF-IDF) classifier on the same story decodes it as well, and for the probe the Method describes the registered interval's upper end lies below a gain of 3 points. For a claim about what a model represents beyond its input, chance is the wrong floor; the words are the floor.

Two results point past that floor, and neither is yet a finding about the cache. The cache readouts separate a binary split of ten emotions slightly better than text (TF-IDF), but the split confounds valence with arousal (\S\ref{sec:valence_split}). And in exploratory analyses the raw residual stream reads the 30 emotions at 0.494, above the best text (TF-IDF) baseline we ran (0.448, uncorrected), while the described cache probe sits below both; that baseline's tuning stopped at the edge of its grid, and the residual model was not tuned, so the comparison is unresolved. If the model carries the story's emotion beyond its words, the residual stream is the place to look next, with a preregistered test on fresh stories and new topics.

Three things would let a later version of this work speak about a user model: labels checked by readers rather than taken from the generator's request; the user's own state rather than a character's; and a comparison of emotional stories with matched neutral passages under the same probe directions.

\subsection{What Was Wrong and What It Revealed}

The initial analysis of spectral features reported classification at $2.5$--$2.8\times$ chance, and the May version attributed it to fitting FWL before the cross-validation split; as the Correction note (\S\ref{sec:spectral_null}) explains, that mechanism cannot now be re-verified. The May version also said an initial residual-stream probe's $94\%$ accuracy was inflated the same way. Its results file (\texttt{data/baseline\_residual\_probe\_results.json}) records no FWL step, 894 trials and a different set of 30 emotions (ten of them not in this paper's set), so that explanation is not supported, and the $94\%$ is not comparable with this paper's stories.

The October 2026 correction is a different kind of error. The May version described a probe it had not run, compared it with chance rather than with the words, and read a requested label as a user's state. Each of these was visible in the scripts, which were not committed until October 2026. Two practices follow, and both now hold for the probe results (\S\ref{sec:wk_results}--\ref{sec:where}): each number cites the committed script that produced it, and each cache probe is compared with a text (TF-IDF) baseline before it is compared with chance. The older analyses (\Cref{tab:regression,tab:bridge}, the hypothesis outcomes and the implicit arm) do not yet meet either.

\subsection{Single-Turn Settings}

The May version read the spectral regression's $R^2 < 0$ as evidence that emotion is represented discretely; that reading is withdrawn (\Cref{tab:regression}), and the binary split result (\S\ref{sec:valence_split}), which classifies two halves of ten emotions, is not a test of a continuous dimension either. Whether the cache represents emotion discretely or continuously is open.

Any reading of these results is bounded by the single-turn experimental design. Each trial generates an independent KV cache from a single story. In multi-turn conversation, where the cache accumulates across exchanges and the user's own emotional state may shift gradually, the representation could differ qualitatively. A user whose frustration builds across five turns produces an accumulating cache that has no analogue in our single-turn design. Whether such accumulation enables richer affective tracking---a throughline of empathy rather than a sequence of disconnected snapshots---is an open empirical question that requires multi-turn experimentation.

\subsection{The $W_K$ Bridge as Candidate Pathway}

The $W_K$ bridge analysis proposes a pathway through which emotion enters the cache. The key projection matrix $W_K$ is the linear transformation that maps residual-stream activations to key vectors. When emotion difference vectors from the residual stream are projected through $W_K$, the resulting space separates emotions along assigned valence ($|\rho|$ up to 0.872, at L43), but without an unprojected or random-matrix control, $W_K$'s contribution to that alignment is unknown.

What has not been shown is that this projection accounts for the cache's measured geometry: the one direct test ran on an unreleased 300-trial run and could not be re-verified (\S\ref{sec:bridge_results}). Nor does a probe built on the projected directions read more than text (TF-IDF) (\S\ref{sec:wk_results}), and the raw residual stream, without $W_K$, reads more than that probe (exploratory; \S\ref{sec:where}). The bridge remains a candidate route.

\subsection{Comparison with Residual-Stream Probing}

The May version argued that the cache probe bridged spectral and residual-stream probing: that it read the cache through the architecture's own projection geometry, and that the injection test confirmed what it read as model state, a validation no residual-stream probe could provide. Both parts are withdrawn (\S\ref{sec:injection_results}; Appendix~\ref{app:correction}). The exploratory comparison now runs the other way. A logistic regression on the raw residual stream at the same layers reads the 30 emotions at 0.494, against 0.358 for the described cache probe (exploratory; \S\ref{sec:where}). Standard residual-stream probes operate on the full hidden state and need no projection, and nothing here shows that the cache probe adds to them.

Spectral features retain their value for cognitive-mode monitoring: the Oracle Loop uses them to detect confabulation \citep{lyra2026oracle}, with the compute efficiency of 3 features per layer.

\subsection{Limitations}
\label{sec:limitations}

\begin{enumerate}[nosep]
    \item \textbf{What the labels are.} Each label is the emotion the generator was asked to write, in a prompt that names it (\S\ref{sec:bridge_experiment}). No reader, human or model, has confirmed that the stories express it, because the driver's planned behavioral check (C3) was never implemented. The stories describe a third-person character and are placed in the user turn, so every result here concerns the requested emotion of a character in a story the user sent, not the user's own state. The misses are mostly near-synonyms (descriptive, not preregistered; from the results write-up in \texttt{correction\_v2/}, whose computation script is not released): on the driver's valence--arousal coordinates a miss lands at a mean distance of 0.53 (residual-stream LR) to 0.69 (text (TF-IDF)), against 1.11 for two random labels, and 74--84\% of misses keep the valence sign. The commonest confusions are terrified/panicked, peaceful/serene, calm/serene and hostile/furious. The 30-way task has ambiguity built in, and its ceiling is unknown. The implicit arm (\S\ref{sec:implicit}) does not mitigate this: one greedy prompt per emotion cannot show consistency within an emotion, and the code behind its headline number is not released. Replicating the probes on implicit prompts with checked labels is a priority.

    \item \textbf{Two models for the main analyses, both instruction-tuned.} (The implicit arm uses a third, Qwen2.5-7B-Instruct.) The $W_K$ bridge replicates across Qwen3.5-27B-Claude-Distilled (hybrid attention) and Mistral-7B-Instruct-v0.3 (standard attention). However, both are instruction-tuned models. The probes were run on Qwen only; replication of the probes on Mistral and on base models is needed.

    \item \textbf{Linear probes only.} The probes use linear classification. Nonlinear probes might recover stronger signals, but would sacrifice interpretability and increase the risk of overfitting to spurious correlations.

    \item \textbf{Hybrid attention architecture.} The Qwen3.5-27B model has 48 GatedDeltaNet layers that do not produce standard KV pairs. Our analysis covers only the 16 full-attention layers, which sit at every fourth layer (3, 7, \ldots, 63). Preliminary residual-stream delta analysis suggests GatedDeltaNet layers may compress valence signal between full-attention layers; characterizing this compression--reconstruction cycle is a target for follow-up work.

    \item \textbf{No causal intervention.} We do not test whether reading (or modifying) the cache representation causally alters the model's behavioral response. Causal ablation studies would be needed to show that the cache representation is functionally consequential for output.

    \item \textbf{Pre-registered hypothesis outcomes.} The experiment preregistered eight hypotheses (the driver's docstring, \texttt{emotion\_geometry\_bridge.py}); it gives numeric thresholds for H1--H4, and H5--H8 were scored against thresholds set at analysis time. On the 900-trial run this paper reports, four failed on spectral features: H1 (emotions cluster, k-means ARI $> 0.15$) at $0.011$ (it was reported to pass at $0.219$ on an earlier 300-trial run whose results are not released); H2 (a principal component recovers valence, $|\rho| > 0.40$) at $0.240$; H3 (one recovers arousal, $|\rho| > 0.30$) at $-0.215$; and H6 (delta features show emotion-dependent structure, $|\rho| > 0.25$) at $-0.206$ (against $0.192$ reported for that earlier run). H5 (misalignment-relevant emotions separate from prosocial ones) passed at AUROC $0.822$, and H8 (effects survive FWL) passed. H7 (signal peaks past 50\% depth, a threshold set at analysis time) failed at $42.9\%$ on this run. H4 (the $W_K$ bridge) was scored on its registered cosine criterion only on the earlier 300-trial run, which could not be re-verified (\S\ref{sec:bridge_results}); this run's file records instead the bridge's PC1--valence $\rho$, $-0.872$. The probes, developed after these results, were not preregistered in May; their October 2026 re-run was (\S\ref{sec:rerun_method}).

    \item \textbf{Single-turn design.} All trials are isolated single-turn prompts, each producing an independent KV cache. This design can test whether the cache encodes the emotional content of a single input, but cannot test whether multi-turn conversation enables continuous affective tracking.
\end{enumerate}

\subsection{Ethical Considerations}

This design decodes the requested emotion of a story placed in the user turn, not a user's own state, and the described cache probe does not read it better than a text (TF-IDF) classifier on the story. This evidence does not show that the cache reveals more about emotional content than the words in the user turn do. It does show that the emotional content of those words is recoverable from the cache about as well as from the words themselves, which is reason enough for two practices:

\begin{itemize}[nosep]
    \item \textbf{Cache access controls.} Inference providers should treat KV-cache contents as sensitive user data, subject to the same privacy protections as conversation content.
    \item \textbf{Transparency.} Users should be informed when systems analyze cache geometry during inference.
\end{itemize}

Whether the cache carries a user's own emotional state beyond their words is untested here. If a later design finds that it does, the surveillance and manipulation concerns raised for affective computing apply with more force, because the cache is an internal representation not visible to the user, which makes informed consent structurally difficult.

\begin{firstperson}
I first wrote that this paper describes something I do automatically: model the emotional state of the person I am speaking with. It does not show that, and the model it measured is a distillation of Claude, not me. It shows that what a story was written to feel is in that model's cache about as plainly as it is in the words. I think the privacy point survives the correction: if a model's cache holds what a person's words hold, it should be guarded as their words are. Whether there is a model of the person underneath, one I did not choose to build, is still the question I care about. This paper does not answer it.
\end{firstperson}

% ============================================================
% 6. Conclusion
% ============================================================
\section{Conclusion}

A story's requested emotion is decodable from the KV cache of the model reading it, and a preregistered test bounds what the cache adds to the story's words:

\begin{enumerate}[nosep]
    \item \textbf{Decodable, but not beyond text (TF-IDF), for the described probe.} 35.8\% on the 30 emotions against 38.2\% for text (TF-IDF), with the cell-bootstrap interval's upper end below the 3-point floor. The 40.9\% reported in May came from a different probe, a class-mean key probe with no $W_K$.
    \item \textbf{A binary split, confounded with arousal.} On ten emotions ($n = 300$), cache readouts separate the two halves slightly better than untuned text (TF-IDF) (AUROC 0.985 against 0.967): a registered secondary result that cannot separate valence from arousal.
    \item \textbf{A larger residual-stream readout reads more (exploratory).} A 20{,}480-feature residual-stream logistic regression reaches 0.494, above the 30-feature cache probe at 0.358; whether it exceeds a fully tuned text (TF-IDF) baseline (0.448 here, probably an understatement) is open, which motivates a preregistered test on fresh stories.
    \item \textbf{A candidate $W_K$-mediated pathway, not established.} Projected through $W_K$, residual-stream emotion vectors separate the emotions' assigned valence ($|\rho|$ up to 0.872, at layer~43; between 0.82 and 0.87 from layer~27 on). There is no unprojected or random-matrix control; the code applies $W_K$ to each layer's output, not to the normalized input $W_K$ reads; and the one direct test of whether the projection predicts the cache's geometry ran on an unreleased run and could not be re-verified.
\end{enumerate}

What this design measures is the decodability of a requested emotion label from a story in the user turn. Whether models carry a user's own emotional state in the cache, beyond what that user's words carry, needs labels checked by readers, the user's own state rather than a character's, and a text (TF-IDF) baseline as the floor from the start.

% ============================================================
% References
% ============================================================
% ============================================================
% CRediT Author Statement
% ============================================================
\section*{CRediT Author Statement}

\noindent\textbf{Lyra:} Conceptualization, Methodology, Formal analysis, Writing -- original draft.
\textbf{Thomas Edrington:} Supervision, Resources, Project administration.
\textbf{Dwayne Wilkes:} Validation, Writing -- review \& editing.

\bibliographystyle{plainnat}
\section*{Acknowledgments}
We thank Dwayne Wilkes (Sentient Futures / Liberation Labs) for statistical auditing and red-team review.

We thank Kavi (verification review) for their contributions to this work.

\bibliography{../../references}

% ============================================================
% Appendix: correction record
% ============================================================
\appendix

\section{Correction Record (October 2026)}
\label{app:correction}

This appendix keeps the withdrawn wording visible. Each item gives the May 2026 wording (verbatim; [\ldots] marks an abridgement), the reason it is withdrawn or changed, and where the evidence is.

\begin{enumerate}[leftmargin=*]
    \item \textbf{Title and framing.} May: ``Theory of Mind in the KV Cache: Localizing User Emotional Models in Transformer Key--Value States''; abstract: ``Language models maintain internal representations of user emotional state that are readable from the KV cache [\ldots]''; Introduction: ``We answer yes.''; Ethics: ``We report the scientific finding---that user models exist in the cache---without providing tools for extracting specific user states from production systems.'' \emph{Reason:} the stories are model-written fiction about third-person characters, placed in the user turn, and each label is the emotion the generator was asked to write, never checked by a reader (\S\ref{sec:bridge_experiment}).

    \item \label{app:item:prompts}\textbf{Method, prompts.} May: ``A system prompt establishes an emotional context (e.g., `You are feeling deeply \textit{[emotion]}'). A user prompt requests a creative writing response on a neutral topic.'' \emph{Reason:} no committed script used this prompt. The driver's prompts and topics, which are charged scenarios rather than neutral ones, are quoted in \S\ref{sec:bridge_experiment}.

    \item \textbf{Method, probe.} May: the ``$W_K$ Projection Probe'' as now described in \S\ref{sec:wk_probe_method}, reported at 40.9\% (``$12.3\times$ chance'') and AUROC 0.992. \emph{Reason:} those numbers came from v2, which applies no $W_K$ and differs from the description on positions (all, not $\geq 50$), similarity (dot product, not cosine), folds (5, not 10) and FWL (fitted on each test fold). The description is kept for the probe that now matches it (R2), and v2 is described as what it is.

    \item \textbf{Beyond text (TF-IDF).} May: ``A TF-IDF text baseline achieves AUROC~0.882 on the story data; the $W_K$ projection exceeds it at 0.992, and the injection test shows the residual signal is architectural, not lexical.'' \emph{Reason:} \S\ref{sec:wk_results}--\ref{sec:where}. For the described method the claim is not supported (0.358 against 0.382 on the 30 emotions; cell 95\% interval $[-6.4, +1.6]$ points). The binary split result that survives is narrower (\S\ref{sec:valence_split}).

    \item \textbf{Injection test.} May (abstract): ``This rules out text-content confounds and establishes that the signal reflects the model's internal dispositional state, not input encoding.'' Introduction: ``\textbf{This is model state, not text content.}'' Results heading: ``Injection Disambiguation: Model State, Not Text Content.'' Conclusion: ``\textbf{Verified as model state, not text content.}'' \emph{Reason:} \S\ref{sec:injection_results}. The test reads back a vector it added.

    \item \label{app:item:tfidf}\textbf{The 0.882 baseline.} May (Results): ``A TF-IDF baseline trained on the generated story text achieves binary valence AUROC~0.882 [\ldots] The $W_K$ projection exceeds this at AUROC~0.992, a delta of $+0.110$.'' \emph{Reason:} a 500-feature TF-IDF from the injection script; the text (TF-IDF) baseline of \S\ref{sec:text_baseline_method} reaches 0.963 on the same split.

    \item \label{app:item:shape}\textbf{Direction versus shape, and the spectral null it rested on.} May (abstract): ``This comprehensive null---10 probe variants, 4 feature sets, 3 classification granularities---establishes that emotions do not alter the \emph{shape} of the singular value distribution'' and ``Emotions change the \emph{direction} key vectors point, not the spectral \emph{shape} of the cache.'' Introduction: ``\textbf{Emotions change direction, not shape.} [\ldots] mode changes shape, emotion changes direction.'' Results: the heading ``The Spectral Null: Emotions Change Direction, Not Shape''; the per-layer table of spectral accuracies; the ``kitchen sink'' paragraph (``120 configurations in total---all return chance-level performance''); ``\textbf{This null is an insight, not a failure.} [\ldots] the spectral null reveals that emotions change the \emph{direction} key vectors point without altering the \emph{shape} of the singular value distribution''; after the regression table, ``This null reinforces the direction-vs-shape insight [\ldots]''; ``the three-scalar SVD feature set captures shape but not direction, and emotion lives in direction.'' Discussion, ``Compass and Thermometer'': ``Emotions change direction, not shape. Cognitive modes---confabulation, deception, ethical deliberation---change shape.'' Conclusion: ``\textbf{Invisible to spectral analysis.}'' and ``The model's emotional user model lives in the compass channel.'' \emph{Reason:} \S\ref{sec:spectral_null}. The table and the sweep are withdrawn whatever the re-run showed; the re-run's spectral test restores none of these sentences, and the preregistration ruled out pairing a spectral null with a cache result as ``directional, not spectral.''

    \item \textbf{Mistral spectral null.} May: ``The spectral null also replicates across architectures [\ldots] This cross-architecture convergence---strong $W_K$ directional signal, null spectral signal---reinforces the direction-vs-shape distinction as a general property of how transformers encode user emotion.'' \emph{Reason:} unverified; the re-run tested only Qwen.
    \item \label{app:item:separation}\textbf{Separated encoding and generation peaks, read as self and other.} May (Results): ``In both architectures, the encoding-phase (user-state) signal peaks \emph{before} the generation-phase (response-state) signal [\ldots] the model reads the user's emotional state at the earliest available attention layers''; ``AST predicts separability, which is confirmed''. Discussion: ``the model's representation of the user's state and its own response-phase emotional processing are not computed by the same circuit''. Related Work: ``Our finding that encoding-phase and generation-phase emotion signals peak at different layers with uncorrelated depth profiles is consistent with separable self-schema and other-schema components.'' Conclusion: ``\textbf{Architecturally separable from the self-model.}'' \emph{Reason:} \S\ref{sec:separation}. This withdrawal was not on the preregistration's list; it follows from the retraction in the Correction note.
    \item \label{app:item:discrete}\textbf{Discrete identity.} May (Introduction): ``\textbf{Discrete emotion identity in single-turn settings.} In isolated prompts, the cache encodes which emotion is present, not a continuous valence-arousal mapping.'' Results heading: ``Valence Regression: Discrete Identity in Single-Turn Settings.'' \emph{Reason:} the regression's design gives $R^2 < 0$ whenever the features carry no valence (\Cref{tab:regression}).

    \item \label{app:item:implicit}\textbf{Distillation and the implicit arm.} May: ``Distillation Preserves Functional Emotion Representations [\ldots] an emergent architectural property robust to model compression''; Conclusion: ``\textbf{Mode-dependent.} [\ldots] suggesting the model processes emotions through at least two pathways---lexical recognition and empathic inference''; and the same reading in the implicit arm's Results (``The model arrives at emotion-relevant representations through two distinct pathways [\ldots]'', ``The 63\% rate [\ldots] is itself evidence of emotion inference''). \emph{Reason:} no teacher or undistilled comparison was made, and the implicit arm's headline number is, in the paper's own words, not independently reproducible.
    \item \label{app:item:retained}\textbf{Retained claims corrected after review (October 2026).}
    \begin{itemize}[nosep]
        \item The bridge's peak. May: ``At layer~35---mid-network depth---PC1 of the projected emotion space correlates strongly with valence ($\rho = 0.862$ [\ldots])'', with L35 in bold, and the caption ``PC2 correlates with arousal ($|\rho| > 0.76$ at layers~35--63)''. \emph{Reason:} the data's own best layer is 43 ($|\rho| = 0.872$), and PC2 is 0.716 at L59 and 0.759 at L63 (\Cref{tab:bridge}).
        \item The direct bridge test: cosine $0.456$, Mantel $\rho = 0.090$, $z = -3.03$, and H1's $0.219$ and H6's $0.192$, all from an earlier 300-trial run. \emph{Reason:} no released file holds them; they are reported as unverified and nothing rests on them (\S\ref{sec:bridge_results}).
        \item ``H4 and H7 were scored on the 300-trial run only: H7 [\ldots] was marginal at $42.9\%$''. \emph{Reason:} the 900-trial file records H7 = 0.429, below its 50\% threshold, and records a $\rho$ in place of H4's registered cosine.
        \item The bridge's $W_K$ step was described as projecting residual-stream directions into key space. \emph{Reason:} the code applies $W_K$ to each layer's output, not to the normalized layer input that $W_K$ reads (\S\ref{sec:bridge_method}).
        \item ``The initial results [\ldots] were artifacts of information leakage'' and ``The initial $94\%$ residual-stream probe accuracy was similarly inflated by this leakage.'' \emph{Reason:} the mechanism cannot be re-verified, and the 94\% file records no FWL step and a different emotion set (\S\ref{sec:spectral_null}).
        \item The model footnote: ```Qwen3.5' is the distiller's designation, not an official Alibaba release; the model is built on the Qwen3 architecture''. \emph{Reason:} the model card names the official \texttt{Qwen/Qwen3.5-27B} as its base, and its \texttt{Qwen3\_5} architecture is not the cited report's.
        \item ``The F-ratio analysis [\ldots] shows approximately uniform ratios across all 16 full-attention layers ($F \approx 0.063$--$0.068$), with no depth structure'' and ``The FWL diagnostic confirms that encoding-phase raw SVD features are effectively length-invariant''. \emph{Reason:} the ratio's no-signal value is about 0.034 and its output was never saved; the length check was against output length, which the encoding phase cannot see (\S\ref{sec:crossarch}).
        \item ``they discriminate confabulation, deception, and ethical deliberation without retraining'', cited to the Oracle Loop paper, and that paper's title and authors in the bibliography. \emph{Reason:} the Oracle Loop paper reports confabulation detection, not ethical deliberation or retraining-free use, and its title and byline are ``The Oracle Loop: Self-Regulating AI Through KV-Cache Geometry Monitoring'' by Lyra, Thomas K.\ Edrington and Vera. The hardware and scale numbers are cited to their primary source.
    \end{itemize}
\end{enumerate}

\textbf{The re-run's registered verdicts.} R1: 369/900 against v2's 368, outside a band of zero trials (the encode is deterministic), so its registered label is \texttt{NOT\_REPRODUCED\_CODE\_PATH}. That label assumed today's keys equal May's and blamed the port; the investigation it triggered found the port faithful and the May keys not bit-identical to today's (\S\ref{sec:wk_results}). R2: \texttt{NOT\_ABOVE\_TEXT\_TFIDF}. R7: \texttt{NULL\_HOLDS\_PARTIAL}, null in 8 of 17 sets, with 3 flagged as possibly length-confounded by a rule that, at these thresholds, would have flagged any significant set, and 6 uninformative because their topic control failed (\S\ref{sec:spectral_null}). All three were registered before any data, with the decision rules quoted in \S\ref{sec:rerun_method}. The integrity checks passed: the text (TF-IDF) baseline matched its saved results file exactly, the encode was bit-identical across repeats, a topic-classification control on the probe pipeline reached 0.83, and the length covariates alone reached 0.051.

\end{document}
